\section{Introduction}\label{sec:intro}

Small clusters bound by pair potentials are a classical object of mathematical chemistry: their equilibria,
saddles and rearrangement pathways organise the energy landscape
\citep{HoarePal1971,HoareMcInnes1983,DoyeWales2002,Trombach2018,SchwerdtfegerWales2024}. For the two-term family
\begin{equation}\label{eq:pot}
\phi(r)=Ar^{-p}-Br^{-q},\qquad A,B>0,\quad p>q>0,
\end{equation}
which contains the Lennard-Jones potential $(p,q)=(12,6)$ and the Mie potentials, symmetric configurations are
equilibria for structural reasons, and their stability can change with the exponents. The eight-particle cube is the
simplest such configuration with a nontrivial spectrum.

A companion work \citep{Paper2} (``Paper~2'') settled the linear problem. With the exact $\Oh$ block decomposition of
the $24$-dimensional displacement space it showed that the cubic equilibrium is a saddle for every $p>q>0$, with Morse
index $2+\one\{q\ge\rstar\text{ or }p>P_A(q)\}+3\cdot\one\{q\ge\qc\text{ or }p>P_T(q)\}\in\{2,3,6\}$, and that zero
modes other than the six rigid motions occur exactly on two disjoint transition curves: $p=P_A(q)$ ($0<q<\rstar$),
where a one-dimensional $\Atwou$ eigenvalue changes sign, and $p=P_T(q)$ ($0<q<\qc$), where the lower eigenvalue of a
$\Ttwog$ block changes sign, with multiplicity three (\cref{fig:plane}). The constants are $\rstar=1.4749107381\ldots$ and
$\qc=2.0199382317\ldots$. An earlier companion preprint \citep{Paper1} (``Paper~1'') studied the constrained radial
problem for the same family.

This paper determines what happens nonlinearly at these curves: which equilibria branch off the cube, on which side
of each curve they exist, what their symmetry and shape are, and how many unstable directions they have. The answer is
local, for each fixed $q$, near the transition point; it rests on equivariant Lyapunov--Schmidt reduction
\citep{GolubitskySchaeffer1985,GolubitskyStewartSchaeffer1988,ChossatLauterbach2000,Field2007} and on two nonlinear
coefficients whose signs we prove with rigorous ball arithmetic.

\paragraph{Results.} Write $\mu$ for the distance $p-P_A(q)$ or $p-P_T(q)$ from the relevant curve, at fixed $q$.
\begin{enumerate}[label=(\roman*)]
\item \emph{The $\Atwou$ pitchfork} (\thmA, \cref{sec:a2u}). The reduced potential is even, its quartic coefficient
is positive along the whole curve (\corM, from the certified \lemM), and two branch equilibria exist for $\mu>0$
only, with amplitude $a^2\propto\mu$. Each consists of two concentric regular tetrahedra of unequal radii on
inversion-related vertex-direction sets and has point group $\Td$.
\item \emph{The $\Ttwog$ representation} (\cref{sec:t2g}). The critical space carries the cubic invariant
$I_3=a_1a_2a_3$; its coefficient $\kappa$ factorises as $qW_q(p-q)\khat$ with $\khat>0$ on the closed curve
(\lemK, certified).
\item \emph{The $\Dthreed$ branches} (\propBone, \cref{sec:d3d}) are transcritical: one branch through the cube with
amplitude $s_*=-(\tau/\kappa)\mu+O(\mu^2)$, on both sides of the curve.
\item \emph{The $\Dtwoh$ branches} (\propBtwo, \cref{sec:d2h}) are $\Z_2$-symmetric pitchforks, supercritical in the
reduced-energy axial convention, with $s_*^2=-(6\tau/c_4)\mu+O(\mu^2)$ for $\mu>0$ only; the sign of the quartic
coefficient $c_4$ is the certified \lemQ. The branch configurations have point group exactly $\Dtwoh$, are not
cuboids, and have rhombic (sheared) cross-sections at first order.
\item \emph{Morse indices of the $\Dtwoh$ branches} (\propBthree\ and \corBthree, \cref{sec:index}). The Hessian of the
reduced potential on the three-dimensional centre space has inertia $(1,0,2)$; the cubic invariant splits the two
transverse directions at order $|s|$. With the frozen block signs of Paper~2 the Hessian on the slice transverse to
the rigid motions has inertia $(4,0,14)$, and in $\R^{24}$ it has inertia $(4,6,14)$, the six zero modes being the
rigid motions. None of these statements is a stability statement.
\end{enumerate}
\Cref{tab:summary} collects the results. Throughout, inertia triples are ordered (negative, zero, positive).

\paragraph{Method and verification.} Every step is either an exact symbolic identity, a standard theorem whose
hypotheses are checked, a frozen result of Paper~2 quoted with its label, or a rigorous interval certificate checked
by an independent implementation (\cref{sec:cert}). Numerical computations are used only for corroboration and are
never proof dependencies; machine-readable dependency graphs make this checkable (\cref{app:graphs}). Each phase was
reproduced in a blind computational audit by a separate AI system, which also re-certified the interval lemmas with
independent implementations (MPFR directed rounding or higher-precision Arb). During the work an implementation defect in the Taylor routines of one
certificate was found; it is documented in \cref{sec:erratum}, together with the evidence that the stored certificate
and the lemma it proves are unaffected.

\paragraph{Relation to the literature.} The equivariant branching lemma \citep{Cicogna1981,GolubitskyStewartSchaeffer1988}
predicts branches along axial isotropy subgroups; for representations with quadratic equivariants, such as the
$\Ttwog$ representation here, Ihrig and Golubitsky \citep{IhrigGolubitsky1984} showed that axial branches are
generically unstable near the bifurcation (see also \citep{ArjevaniField2022}). That is a stability statement; whether
a given axial branch is transcritical or a pitchfork is decided axis by axis, and we derive both types directly from
the implicit function theorem. Bushes of nonlinear vibrational modes in octahedral Lennard-Jones systems were studied
in \citep{Chechin2003}. We are not aware of a prior complete classification of these local branches for this
specific eight-particle two-term inverse-power system.

\paragraph{Organisation.} \Cref{sec:model,sec:symmetry,sec:curves} fix the model, the symmetry conventions and the
frozen inputs from Paper~2. \Cref{sec:a2u} treats the $\Atwou$ curve, \cref{sec:t2g,sec:d3d,sec:d2h,sec:index} the
$\Ttwog$ curve. \Cref{sec:cert} describes the certificates, the audits and the erratum; \cref{sec:discussion} discusses
limitations, and \cref{sec:repro} reproducibility.
